\section{Classical preliminaries}  \label{sec:prelims}

\subsection{Finite fields and polynomials}
In this section we review some basic facts about finite fields and
polynomials over them. These facts can be found in a standard
reference such as~\cite{MBG+13}. Let $p$ be a prime and $q = p^t$ be a prime power. We denote by $\F_p$
and $\F_q$ the finite fields with $p$ and $q$ elements,
respectively. The field $\F_p$ is called the \emph{base field} or
\emph{prime subfield} of $\F_q$. The larger field $\F_q$ can be viewed as a
$t$-dimensional vector space $\F_p^{t}$ over the smaller field.
We define the trace $\tr: \F_q \to \F_p$ by
\[ \tr[a]  = \sum_{\ell=0}^{t-1} a^{p^\ell}. \]
The trace is a linear map under linear combinations with coefficients
drawn from $\F_p$. 

A \emph{basis} for
$\F_q$ over $\F_p$ consists of $k$ elements $\{\alpha_1, \dots,
\alpha_k\}$, such that any element $u \in \F_q$ can be written as a
linear combination
\[ u = \sum_{i=1}^k c_i \alpha_i, \]
where the coefficients $c_i \in \F_p$. Two bases $\{\alpha_i\}$ and
$\{\beta_i\}$ are \emph{dual} bases if $\tr[\alpha_i \beta_j]
= \delta_{ij}$, where $\delta_{ij}$ is the Kronecker delta function.

\paragraph{Fields of characteristic $2$}:
When $p =2$ (i.e. $q$ is even), several useful properties hold. Most
importantly for us, the field $\F_q$ has a \emph{self-dual} basis: that is, there exists a basis $\{\alpha_i\}$ such
that $\tr[\alpha_i \alpha_j] =
\delta_{ij}$~\cite[Theorem 1.9]{MBG+13}. This means that
given a field element $u = \sum_i c_i \alpha_i$, we can recover the
coefficient $c_j$ by the expression $c_j = \tr[u \alpha_j]$.

\paragraph{Fourier identities}
Below, we give two useful identities for simplifying Fourier sums over
finite fields. We set $\omega = e^{2\pi i / p}$ to be a $p$-th root of unity.

\begin{fact}\label{fact:averages-to-zero}
$\E_{\bu \in \F_q} \omega^{\tr[\bu \cdot a]} = 0$ if $a$ is nonzero.
\end{fact}
\begin{proof}
  If $a \neq 0$, then there must exist some nonzero $y \in \F_q$ such that
  $\tr[ay] \neq 0$. Let the value of the expectation we want to compute be
  denoted by $\sigma$. Then we have
  \begin{align*}
    \sigma &= \E_{\bu \in \F_q} \omega^{\tr[a \bu]} \\
           &= \E_{\bu \in \F_q} \omega^{\tr[a(\bu + y)]} \\
           &= \omega^{\tr[ay]} \E_{\bu \in \F_q}\omega^{\tr[a\bu]} \\
           &= \omega^{\tr[ay]} \sigma,
  \end{align*}
    and thus $\sigma = 0$.
  \end{proof}

  \begin{fact}\label{fact:averages-to-zero-subspace}
    Let $V$ be a subspace of $\F_q^n$. Then $\E_{\bu \sim V}
    \omega^{\tr[\langle \bu, a \rangle]} = 0$ if $a \notin V^\perp$.
  \end{fact}
  \begin{proof}
    The idea is the same as the proof of the previous fact. Suppose $a
    \notin V^\perp$. Then there exists some nonzero $y \in \F_q^n$
    such that $\tr[\langle a, y \rangle] \neq 0$. Letting the value of
    the expectation we wish to compute be denoted $\sigma$, we have
    \begin{align*}
      \sigma &= \E_{\bu \sim V} \omega^{\tr[\langle \bu, a \rangle]}
      \\
             &= \E_{\bu \sim V} \omega^{\tr[ \langle \bu + y,
               a\rangle]} \\
             &= \omega^{\tr[\langle a, y\rangle]} \sigma,
    \end{align*}
    and hence $\sigma $ must vanish.
  \end{proof}
  


\subsection{Two-player one-round games and $\MIP$}
A two-prover nonlocal game is an interaction between a verifier and
two noncommunicating provers, in which the verifier samples a pair of random
questions and sends them to the provers, receives a pair of answers,
and decides whether to accept or reject based on the questions and answers. In the literature, a game is usually
taken to be described by the verifier's distribution over question
pairs, together with a table describing the verifier's behavior
for all possible choices of questions and answers. For our purposes,
it will be more convenient to work with \emph{uniformly generated} families of
games, which are specified by Turing machines that sample the
questions and decide whether to accept or reject given the questions
and answers.
\begin{definition}[Two-player one-round uniform game family]\label{def:two-player-one-round}
A \emph{two-prover one-round game uniform game family} $\game$
is an interaction between a verifier and two provers, Alice and Bob.
The verifier $V = (\mathrm{Alg_{\mathrm{Q}}}, \mathrm{Alg_{\mathrm{A}}})$
consists of a ``question" randomized Turing machine $\mathrm{Alg_{\mathrm{Q}}}$
and an ``answer" deterministic Turing machine $\mathrm{Alg_{\mathrm{A}}}$.
Given an input string $\mathsf{input}$, the verifier samples two questions
$(\bx_0, \bx_1) \sim \mathrm{Alg_{\mathrm{Q}}}(\mathsf{input})$ and distributes~$\bx_0$ to Alice and~$\bx_1$ to Bob.
They reply with answers~$\ba_0$ and~$\ba_1$, respectively, and the verifier accepts if
$\mathrm{Alg_{\mathrm{A}}}(\mathsf{input}, \bx_0, \bx_1, \ba_0, \ba_1) = 1$.
A strategy for Alice and Bob is said to be \emph{classical} if they are allowed shared randomness but no shared quantum resources.
The \emph{value} of Alice and Bob's strategy is simply the probability that the verifier accepts,
and the \emph{classical value} of the game is the maximum value of any classical strategy.
We write $\qlength{\game}$ for the maximum bit length
of the questions as a function of the input $\mathsf{input}$, and similarly
$\alength{\game}$ for the maximum bit length of the answers,
$\qtime{\game}$ for the maximum running time of $\mathrm{Alg}_{\mathrm{Q}}$,
and $\atime{\game}$ for the maximum running time of
$\mathrm{Alg}_{\mathrm{A}}$. Often we will not explicitly write the
dependence of these quantities on $\mathsf{input}$.
\end{definition}

\begin{definition}[Multiprover interactive proofs]\label{def:mip-protocols}
A \emph{2-player 1-round multiprover interactive proof protocol} is a uniform
game family~$\game$ as in \Cref{def:two-player-one-round}. For
parameters $0 < s < c \leq 1$, we say that
the protocol $\game$ decides the language $L$ with completeness
$c$ and soundness $s$ if the following three conditions are true.
\begin{itemize}
\item[$\circ$] (Completeness) Suppose $\mathsf{input} \in L$.  Then there is a classical strategy for~$\game$ with value at least~$c$.
\item[$\circ$] (Soundness) Suppose $\mathsf{input} \notin L$.  Then every classical strategy for~$\game$ has value at most~$s$.
\item[$\circ$] All of $\qlength{\game}$, $\alength{\game}$,
  $\qtime{\game}$, and $\atime{\game}$ are $\poly(n)$ where $n$ is
  the bit length of $\mathsf{input}$.
\end{itemize}
The class $\MIP_{c, s}$ is the set of all languages that can be
decided by multiprover interactive proof protocols with the parameters $c,s$.

If $c - s$ is a constant, then we will suppress the dependence on them when writing $\MIP$ and just say that $L \in \MIP$.
Here, ``$c$" is referred to as the \emph{completeness} and ``$s$" is referred to as the \emph{soundness}.
We will typically deal with the case when $c = 1$ and $s = 1-\eps$, where $\eps > 0$ is a small constant.

\end{definition}
In this definition of $\MIP$, the parameters $\qlength{\game}, \alength{\game},
\qtime{\game}, \atime{\game}$ are required to be polynomial in the
input length $n$. However, in this paper, several of the intermediate
results we achieve are protocols where these parameters scale
superpolynomially (indeed, even exponentially or worse) in $n$. In
these cases, we will explicitly indicate the dependence of these
parameters on $n$.


\subsection{Low-degree code}

Let $q$ be a prime power and $h\leq q$ be an integer.  
Let $H$ be a subset of $\F_q$ of size~$h$.
For $n \geq 0$, let $x \in H^n$.
The \emph{indicator function of~$x$ over $H^n$}
is the polynomial with inputs $y \in \F_q^m$ defined as
\begin{equation*}
\indicator{H}{x}{y} := \frac{\prod_{i=1}^m \prod_{b \in H, b \neq x_i}
  (y_i - b)}{\prod_{i=1}^m \prod_{b \in H, b \neq x_i} (x_i - b)}.
\end{equation*}
There are two properties of this polynomial that we will need:
\begin{enumerate}
\itemsep -.5pt
\item[(i)] that it is low-degree, i.e. a degree-$m(h-1)$ polynomial,
\item[(ii)] that for any $x, y \in H^m$, $\indicator{H}{x}{y} = 1$ if and only if $x = y$,
and otherwise $\indicator{H}{x}{y}= 0$.
\end{enumerate}
Using this, we can define the low-degree code.

\begin{definition}[Low-degree encoding]
Let $|\calS| \leq h^m$, and let $\pi:\calS \rightarrow H^m$ be an injection.
Then the \emph{low-degree encoding} (sometimes also called the \emph{Reed-Muller encoding})
of a string $a \in \F_q^{\calS}$ is the polynomial $g_a :\F_q^m \rightarrow \F_q$ defined as
\begin{equation*}
g_a(x) := \sum_{i\in\calS} a_i \cdot \indicator{H}{\pi(i)}{x}.
\end{equation*}
\end{definition}
By the properties of the indicator function above,
(i) $g_a$ is a degree-$m(h-1)$ polynomial,
and (ii)~$g_a(\pi(i)) = a_i$ for all $i \in \calS$.
We will typically, though not always, take $\calS = [n]$.
Given an error-correcting code, there are two key properties we care about:
the rate and the distance.
The rate of the low-degree code is $n / q^m$.
As for the distance, we can estimate it with the following lemma.

\begin{lemma}[Schwartz-Zippel lemma~\cite{Sch80,Zip79}]\label{lem:schwartz-zippel}
Let $f, g$ be two unequal $m$-variate degree-$d$ polynomials over $\F_q$. Then
\begin{equation*}
\Pr_{\bx \sim \F_q^m}[f(\bx) = g(\bx)] \leq d/q.
\end{equation*}
\end{lemma}
As a result, the low-degree encoding has relative distance $m(h-1) / q$.
In a typical application, we would like a code with large rate and distance.
To achieve this, we will often use the following ``rule of thumb" setting of parameters:
\begin{equation}\label{eq:parameters}
h = \Theta(\log(n)),
\qquad
m = \Theta\left(\frac{\log(n)}{\log\log(n)}\right),
\qquad
q = \mathrm{polylog}(n).
\end{equation}
This gives a code with rate $1/\mathrm{poly}(n)$ and distance $o(1)$.
The polynomials involved are degree $d = \Theta(\log(n)^2/\log\log(n))$.

\subsection{A canonical low-degree encoding}

The low-degree encoding affords us some flexibility when choosing the parameters and the injection;
however, for our application we will have to choose these with care,
because each of our uses of the low-degree code requires that the injection~$\pi$ be efficiently computable.
In this section, we give a simple, canonical choice for the subset~$H$ and the injection~$\pi$ so that this is true.

\begin{definition}\label{def:admissible-prelims}
We say that
$n$,
$h= 2^{t_1}$,
$q = 2^{t_2}$,
and $m$
are \emph{admissible parameters} if $t_1 \leq t_2$ and $h^m\geq n$.
\end{definition}

The following definition gives the canonical encoding.

\begin{definition}[Canonical low-degree encoding]\label{def:canonical-low-degree}
Let $n$, $h = 2^{t_1}$, $q= 2^{t_2}$, and $m$ be admissible parameters.
Set $\ell = t_1 \cdot m$.
The \emph{canonical low-degree code} is defined as follows.
\begin{itemize}
\item[(i)] Let $e_1, \ldots, e_{t_2}$ be a self-dual basis for~$\F_q$ over~$\F_2$. Then we set $H$ to be the subset
	\begin{equation*}
		H:= H_{t_1, t_2} = \{b_1 \cdot e_1 + \cdots + b_{t_1} \cdot e_{t_1} \mid b_1, \ldots, b_{t_1} \in \F_2\}.
	\end{equation*}
	As desired, $|H| = h$.
\item[(ii)] Let $\sigma:= \sigma_{t_1, t_2}: \{0, 1\}^{t_1} \rightarrow H_{t_1, t_2}$ be the bijection $\sigma(b_1, \ldots, b_{t_1}) = b_1 \cdot e_1 + \cdots + b_{t_1} \cdot e_{t_1}$.
	From this, we can construct a bijection $\sigma_{\ell, t_1, t_2} : \{0, 1\}^\ell \rightarrow H^m$ by setting
	\begin{equation*}
		\sigma_{\ell, t_1, t_2}(b_1, \ldots, b_\ell) = (\sigma(b_1, \ldots, b_{t_1}), \sigma(b_{t_1+1}, \ldots, b_{2t}), \ldots, \sigma(b_{\ell-t_1+1}, \ldots, b_{\ell})).
	\end{equation*}
\item[(iii)] Given an index $i \in [n]$, write $\mathrm{bin}_\ell(i)$ for its $\ell$-digit binary encoding. 
	Then we define the injection $\pi:=\pi_{\ell, t_1, t_2}:[n] \rightarrow H^m$ as $\pi(i) = \sigma_{\ell, t_1, t_2}(\mathrm{bin}_\ell(i))$.
\end{itemize}
\end{definition}

The following proposition gives the time complexity of the canonical low-degree encoding.

\begin{proposition}\label{prop:canonical-time}
The bijection~$\sigma_{\ell, t_1, t_2}$ and the injection~$\pi:=\pi_{\ell, t_1, t_2}$ are both computable in time $m \cdot \mathrm{poly log}(q)$.
As a result, given a string $a \in \F_q^n$ and a point $x \in \F_q^m$, the value $g_a(x)$ takes time $\poly(n, m, q)$ to compute.
\end{proposition}

\subsection{Low-degree testing}

\begin{definition}[Surface-versus-point test]
The \emph{surface-versus-point low-degree test with parameters} $m$, $d$, $q$ (a prime power), and~$k$, denoted $\game_{\mathrm{Surface}}(m, d, q, k)$, is defined as follows.
Let $\bv_1, \ldots,\bv_k$ be~$k$ uniformly random vectors in $\F_q^m$, and
let $\bs$ be a uniformly random affine subspace parallel to
$\mathrm{span}\{\bv_1, \ldots, \bv_k\}$ (that is, $\bs$ is the set $\{\bw + \lambda_1 \bv_1 +\cdots
+\lambda_k \bv_k : \lambda_1, \ldots, \lambda_k \in \F_q\}$ for a uniformly
random $\bw$),
and let $\bu$ be a uniformly random point on~$\bs$.
Given these, the test is performed as follows.
\begin{itemize}
\itemsep -.5pt
\item[$\circ$] The vectors $\bv_1, \ldots, \bv_k$ and the surface $\bs$ are given to Alice, who responds with a degree-$d$ polynomial $\boldf: \bs \rightarrow \F_q$.
\item[$\circ$] The point $\bu$ is given to Bob, who responds with a number $\bb \in \F_q$.
\end{itemize}
Alice and Bob pass the test if $\boldf(\bu) = \bb$.
\end{definition}

\begin{remark}
Let us remark briefly on the encodings used in this test.
A surface~$s$ with directions $v_1, \ldots, v_k$
is encoded by the string $(u, w_1, \ldots, w_k) \in \F_q^{(k+1) n}$.
Here, $u$ is the lexicographically minimum point in~$s$,
and $w_1, \ldots, w_k$ are the rows of the matrix produced by taking the matrix with rows $v_1, \ldots, v_k$
and transforming it to reduced row echelon form.
We note that given $v_1, \ldots, v_k$ and a point $u \in s$,
this encoding can be produced in time $\poly(n, k, \log(q))$.

A function $f:s\rightarrow \F_q$ is \emph{a degree-$d$ polynomial on~$s$}
if there exists a degree-$d$ $k$-variate polynomial $f':\F_q^k \rightarrow \F_q$
such that $f'(\lambda_1, \ldots, \lambda_k) = f(u + \lambda_1 w_1 + \cdots + \lambda_k w_k)$.
When~$s$ is already known, we can encode~$f$ by specifying~$f'$,
which involves writing out its $d[k] := \binom{d+k}{k}$ coefficients in some arbitrary but fixed order.
\end{remark}

We note that this definition of the surface-versus-point test differs slightly from the standard
definition of the surface-versus-point test in two respects. First, we
do not require that the vectors $\bv_1, \ldots, \bv_k$ be linearly independent
or even nonzero, which implies that there is a
\begin{equation*}
1-
\left(1 -\frac{1}{q^m}\right)
\left(1 -\frac{q}{q^m}\right)
\cdots
\left(1 -\frac{q^{k-1}}{q^m}\right)
\leq \frac{q^k}{q^m}
\end{equation*}
 chance that $\bs$
is less than $k$-dimensional. Second, we send the vectors $\bv_1,
\ldots, \bv_k$ to Alice in addition to the description of the surface $\bs$. It
is not hard to see that these two modifications do not asymptotically harm the
soundness guarantee obtained for the standard plane-versus-point test shown by Raz and
Safra~\cite{RS97}, which we restate here.

\begin{theorem}[\cite{RS97}]\label{thm:raz-safra}
There exist absolute constants $c, c' > 0$ such that the following holds.
Suppose Alice and Bob pass $\game_{\mathrm{Surface}}(m,d,q,2)$ with probability at least $\mu$.
Then there exists a degree-$d$ polynomial $g:\F_q^m \rightarrow \F_q$ such that
\begin{equation*}
\Pr_{(\bs, \bu)}[g(\bu) = \bb] \geq \mu - c\cdot m (d/q)^{c'}.
\end{equation*}
\ignore{where $\bb$ is Bob's random answer given the point query~$\bu$.}
\end{theorem}

Explicit values for $c, c'$ have been derived by Moshkovitz and Raz~\cite{MR08},
albeit for the weaker guarantee that~$g$ be a degree-$m d$ polynomial,
which is still sufficient for most applications.

\paragraph{Communication cost.}
We can compute the communication cost of this test as follows.

\begin{itemize}
\itemsep -.5pt
\item[$\circ$] \textbf{Question length:} We encode a plane in $\F_q^m$ with a string $(u, v_1, v_2) \in \F_q^{3m}$. This requires $3 m \log(q)$ bits to communicate.
\item[$\circ$] \textbf{Answer length:} A degree-$d$ bivariate polynomial on~$\F_q$ can be described with $\binom{d+2}{2} \leq (d+1)^2$ coefficients in~$\F_q$. These require $(d+1)^2 \log(q)$ bits to communicate.
\end{itemize}
\noindent
Recalling \Cref{eq:parameters}, a typical setting of parameters
gives questions of length $\Theta(\log(n))$,
and  answers of length $\Theta(\log(n)^4/\log\log(n))$.

\subsection{Simultaneous low-degree testing}\label{sec:simultaneous-classical}

\begin{definition}[Simultaneous surface-versus-point test]\label{def:simultaneous-plane-v-point}
The \emph{simultaneous surface-versus-point low-degree test with parameters} $m$, $d$, $q$ (a prime power), $k$, and $\ell$, denoted $\game_{\mathrm{Surface}}^\ell(m, d, q, k)$,
 is defined as follows.
A draw $(\bs, \bu)$ is sampled as in $\game_{\mathrm{Surface}}(m,d,q,k)$. Given this, the test is performed as follows.
\begin{itemize}
\itemsep -.5pt
\item[$\circ$] The surface $\bs$ is given to Alice, who responds with $\ell$ degree-$d$ polynomials $\boldf_1, \ldots, \boldf_{\ell}: \bs \rightarrow \F_q$.
\item[$\circ$] The point $\bu$ is given to Bob, who responds with $\ell$ numbers $\bb_1, \ldots, \bb_{\ell} \in \F_q$.
\end{itemize}
Alice and Bob pass the test if $\boldf_1(\bu) = \bb_1, \ldots,
\boldf_{\ell}(\bu) = \bb_{\ell}$.
\label{def:simul}
\end{definition}
Classically, the $k=2$ case of this test can be reduced to a slight generalization of \Cref{thm:raz-safra} using a
simple and standard union-bound argument.
Quantumly, however, a corresponding entanglement-sound analogue of this generalization is not known to hold.
Instead, we use a slightly more involved reduction in which the~$\ell$ outputs of Alice and Bob are ``combined" to create a strategy for the ``standard" plane-versus-point test.
(This technique is standard and was also used in the proof of Lemma~4.6 in~\cite{NV18a} for the case $\ell =2$.)
In this section, we will introduce the notation needed for this reduction
and carry out the proof of the classical soundness of the simultaneous low-degree test
as a warm-up for our proof of quantum soundness later.
We begin by showing how to combine~$\ell$ functions by introducing~$\ell$ ``indexing" variables.

\begin{notation}
Let $g_1, \ldots, g_\ell:s \rightarrow \F_q$ be functions, where~$s$ is a subset of $\F_q^m$.
Then we define the new function $\mathrm{combine}_g(x, y): \F_q^\ell \otimes s \rightarrow \F_q$ as follows:
\begin{equation*}
\mathrm{combine}_g(x, y) = x_1\cdot g_1(y) + \cdots + x_\ell \cdot g_\ell(y).
\end{equation*}
We will typically apply this with $s = \F_q^m$ or~$s$ a dimension-$k$ subspace of~$\F_q^m$.
\end{notation}

If the $g_i$'s are degree-$d$ polynomials on $\F_q^m$, this produces a degree-$(d+1)$ polynomial on $\F_q^{\ell + m}$.
First, we show that given a surface-versus-point query from this $(\ell + m)$-dimensional space,
we can produce a surface-versus-point query from the $m$-dimensional space.

\begin{proposition}\label{prop:really-the-same-dist}
Given a subset $s \subseteq \F_q^{\ell + m}$, let $s_{\mathrm{proj}} = \{y \mid (x, y) \in \F_q^{\ell} \otimes \F_q^m\}$.
\begin{itemize}
\item[$\circ$] If~$s$ is a dimension-$k$ subspace of $\F_q^{\ell+m}$, then $s_{\mathrm{proj}}$ is a dimension-$k'$ subspace of $\F_q^{m}$, for $k' \leq k$.
\end{itemize}
Define $\bs' \sim_k s_{\mathrm{proj}}$ to be a uniformly random dimension-$k$ subspace of $\F_q^{m}$ containing $s_{\mathrm{proj}}$.
\begin{itemize}
\item[$\circ$] If $\bs$ and $(\bx, \by)$ are distributed as $\mathscr{D}_{\mathrm{Surface}}(\ell + m,q,k)$,
		then $\bs' \sim \bs_{\mathrm{proj}}$ and $\by$ are distributed as $\mathscr{D}_{\mathrm{Surface}}(m, q, k)$.
\end{itemize}
\end{proposition}
\begin{proof}
The first bullet follows because if $\{(x_1, y_1), \ldots, (x_k, y_k)\}$ is a set of~$k$ linearly independent vectors which span~$s$,
then $\{y_1, \ldots, y_k\}$ is a set of~$k$ vectors which span~$s_{\mathrm{span}}$, though they may no longer be linearly independent.
The second bullet follows by symmetry.
\end{proof}

Next, we show that answers to the queries on the $m$-dimensional space
can be used to produce answers to the queries on the $(\ell+m)$-dimensional space.

\begin{proposition}\label{prop:sub-subspace}
Let $s$ be a dimension-$k$ subspace of $\F_q^{\ell + m}$, and let $s' \subseteq \F_q^{m}$ be a subspace which contains $s_{\mathrm{proj}}$.
Then $\F_q^\ell \otimes s'$ is a subspace, and it contains~$s$.
In particular, if $f_1, \ldots, f_\ell$ are degree-$d$ functions on $s'$, then $\mathrm{combine}_f$ is a degree-$(d+1)$ function on $\F_q^\ell \otimes s'$,
and it can be restricted to a degree-$(d+1)$ function on~$s$.
\end{proposition}
\begin{proof}
Consider a point $(x, y) \in s$. Then $y \in s_{\mathrm{proj}} \subseteq s'$, and so $(x, y) \in \F_q^{\ell} \otimes s$.
The statement about $\mathrm{combine}_f$ follows immediately.
\end{proof}

Finally, we need a technical result: that nonlinear low-degree polynomials rarely become linear after restricting variables.

\begin{definition}
Let $n \geq 0$.
A function $f:\F_q^{\ell + n} \rightarrow \F_q$ is \emph{exactly linear in~$x$}
if it can be written as
\begin{equation*}
f(x, y) = x_1 \cdot f_1(y) + \cdot + x_\ell \cdot f_{\ell}(y).
\end{equation*}
(We do not allow constant terms.)
Note that when $n = 0$, such a function can be written as $c_1 \cdot x_1 + \cdots + c_\ell \cdot x_\ell$, where each $c_i \in \F_q$,
in which case we simply call it ``exactly linear".
Given a function $f(x, y)$ and a string $y \in \F_q^m$,
we will also write $f|_y$ for the function defined as $f_y(x) = f(x, y)$.
\end{definition}

\begin{proposition}\label{prop:exactly-linear-prop}
Suppose $f(x, y): \F_q^{\ell + m} \rightarrow \F_q$
is a degree-$d$ polynomial which is not exactly linear in~$x$.
Then the probability that $f|_{\by}$ is exactly linear, over a uniformly random $\by \sim \F_q^m$, is at most $d/q$.
\end{proposition}
\begin{proof}
Because~$f$ is not exactly linear in~$x$, it contains some non-linear $x$-monomial $x^i = x_1^{i_1} \cdots x_\ell^{i_\ell}$
in which $i_1 + \cdots + i_\ell$ is either zero or at least two.
Thus, $f$ can be written as
$f(x, y) = x^i \cdot g_i(y) + f'(x, y)$, where $g_i(y)$ is degree-$d$ and $f'$ contains no~$x^i$ terms.
For $f|_{\by}$ to be exactly linear, this term must vanish, which means $g_i(\by) = 0$.
But by Schwartz-Zippel (\Cref{lem:schwartz-zippel}), this happens with probability at most $d/q$.
\end{proof}

We are now ready to prove soundness of the simultaneous low-degree test in the $k = 2$ case.

\begin{theorem}\label{thm:simultaneous-raz-safra}
There exists absolute constants $c, c' >0$ such that the following holds.
Suppose Alice and Bob pass $\game_{\mathrm{Surface}}^\ell(m, d, q,2)$ with probability at least $\mu$.
Then there exist degree-$d$ polynomials $g_1, \ldots, g_{\ell} : \F_q^m \rightarrow \F_q$ such that
\begin{equation*}
\Pr_{(\bs, \bu)} [g_1(\bu) = \bb_1, \ldots, g_{\ell}(\bu) = \bb_\ell] \geq \mu - c \cdot (m + \ell) (d/q)^{c'}.
\end{equation*}
\end{theorem}

\begin{proof}
Let $c, c' > 0$ be as in \Cref{thm:raz-safra}.
We pick the constants in this theorem, say $\hat{c}, \hat{c}'$  so that
\begin{equation*}
\mu - c \cdot (m + \ell) ((d+1)/q)^{c'} - 2(d+1)/q \geq \mu - \hat{c} \cdot (m + \ell) (d/q)^{\hat{c}'}.
\end{equation*}
Note that this means that the theorem is trivial when $2(d+1)/q \geq  \mu - c \cdot (m + \ell) ((d+1)/q)^{c'}$.
As such, we will assume below that
\begin{equation}\label{eq:really-technical}
2(d+1)/q <  \mu - c \cdot (m + \ell) ((d+1)/q)^{c'}.
\end{equation}

Suppose Alice and Bob pass $\game_{\mathrm{Surface}}^\ell(m, d, q,2)$ with probability at least $\mu$.
We will use them to simulate two provers, ``Combined Alice" and ``Combined Bob", who pass
the single-function low-degree test $\game_{\mathrm{Surface}}(\ell + m, d+1, q,2)$ with probability at least $\mu$.
They are specified as follows:
\begin{itemize}
\item[$\circ$] \textbf{Combined Alice:} Given $\bs \subseteq \F_q^{\ell + m}$, draw $\bs' \sim_2 \bs_{\mathrm{proj}}$.  Give it to Alice, who responds with $\boldf_1, \ldots, \boldf_{\ell}: \bs' \rightarrow \F_q$.  Output the function $\mathrm{combine}_{\boldf}|_{\bs}$.
\item[$\circ$] \textbf{Combined Bob:} Given $(\bx, \by) \in \F_q^{\ell + m}$, compute $\by \in \F_q^m$. Give it to Bob, who responds with $\bb_1, \ldots, \bb_\ell \in \F_q$.
			Return $\mathrm{combine}_{\bb}(\bx) \in \F_q$.
\end{itemize}
By \Cref{prop:really-the-same-dist}, $\bs'$ and $\by$ are distributed as the questions in $\game_{\mathrm{Surface}}^\ell(m, d, q,2)$.
Using our assumption on Alice and Bob, this means that $\boldf_1(\by) = \bb_1$, \ldots, $\boldf_\ell(\by) = \bb_\ell$ with probability at least $\mu$.
As a result, $(\mathrm{combine}_{\boldf}|_{\bs})(\bx, \by) = \mathrm{combine}_{\bb}(\by)$ with probability at least~$\mu$.
By \Cref{prop:sub-subspace}, $\mathrm{combine}_{\boldf}|_{\bs}$ is a degree-$(d+1)$ function on~$\bs$,
and so it is a valid response to subspace queries.
This means Combined Alice and Bob pass~$\game_{\mathrm{Surface}}(\ell + m, d+1, q,2)$ with probability at least~$\mu$.

Thus, we can apply \Cref{thm:raz-safra}. It gives a degree-$(d+1)$ function $g: \F_q^{\ell + m} \rightarrow \F_q$ such that
\begin{equation}\label{eq:g-equals-combine}
\Pr_{\bx, \by}[g(\bx, \by) = \mathrm{combine}_{\bb}(\bx)] \geq \mu - c \cdot (\ell + m)\cdot ((d+1)/q)^{c'}.
\end{equation}
We would like to show that~$g$ is exactly linear in~$x$.
Assume for the sake of contradiction that this is not the case.
Because $\bb$ depends only on~$\by$ (and Bob's internal randomness),
we can consider varying these two variables independently of~$\bx$.
By \Cref{prop:exactly-linear-prop}, the probability that $g_{\by}$ is not exactly linear is at least $1-(d+1)/q$.
In this case, because $\mathrm{combine}_b(\bx)$ is always exactly linear,
the probability that $g|_{\by}(\bx) = \mathrm{combine}_{\bb}(\bx)$ is at most $(d+1)/q$ by Schwartz-Zippel (\Cref{lem:schwartz-zippel}).
As a result, the probability that $g(\bx, \by) = \mathrm{combine}_{\bb}(\bx)$ is at most $(d+1)/q + (d+1)/q$, which contradicts \Cref{eq:really-technical,eq:g-equals-combine}.
Thus, we may conclude that~$g$ is exactly linear in~$x$.

This implies that we can write $g(x, y) = \sum_i x_i \cdot g_i(y)$, where each $g_i$ is a degree-$d$ polynomial.
Now, for any fixed~$b$ and~$y$, if it is not the case that $g_1(y) = b_1$, \ldots, $g_\ell(y) = b_\ell$,
then the probability that $g(\bx, y) = \mathrm{combine}_{b}(\bx)$ over a random~$\bx$ is at most $1/q$ by Schwartz-Zippel since both are exactly linear functions.
Thus, if $\eta$ is the probability that $g_1(\by) = \bb_1$, \ldots, $g_\ell(\by) = \bb_\ell$,
then the probability that $g(\bx, \by) = \mathrm{combine}_{\bb}(\bx)$ is at most $\eta + (1-\eta)/q \leq \eta + 1/q$.
Combined with \Cref{eq:g-equals-combine}, this implies the theorem.
\end{proof}

\subsection{$\NEXP$, $\NEEXP$, and complete problems for them}

\begin{definition}
  \label{def:neexp}
  The class $\neexp$ (respectively, $\NEEXP$) is the class of all problems that can be solved
  in exponential (respectively, doubly-exponential) time by a nondeterministic Turing
  machine. Formally,
\begin{equation*}
\NEXP = \bigcup_{c \in \mathbb{N}}  \NTIME(2^{n^c}),\qquad
\NEEXP = \bigcup_{c \in \mathbb{N}}  \NTIME(2^{2^{n^c}}).
\end{equation*}
\end{definition}

A standard way of generating $\NEXP$-complete problems 
is by considering ``succinct" versions of $\NP$-complete problems,
in which an exponential-sized input is encoded by a polynomial-sized circuit.
The canonical complete problem is a succinct version of $\mathsf{3Sat}$,
but there is considerable freedom in choosing the succinct encoding used.
We choose the following encoding.

\begin{definition}
$\succinct$ is the following problem.
\begin{enumerate}
\item[$\circ$] \textbf{Input:} a circuit~$\calC$ with $3n+3$ input bits and size $\mathrm{poly}(n)$.
It encodes the 3-Sat instance $\psi_{\calC}$ with variable set $x_u$ for $u \in \{0, 1\}^n$
which includes the constraint $(x_{u_1}^{b_1} \lor x_{u_2}^{b_2}\lor x_{u_3}^{b_3})$ whenever
\begin{equation*}
\calC(u_1, u_2, u_3, b_1, b_2, b_3) = 1.
\end{equation*}
(Here, $x_i^1$ refers to the literal $x_i$ and $x_i^0$ refers to the negated literal $\overline{x_i}$.)
\item[$\circ$] \textbf{Output:} accept if $\psi_{\calC}$ is satisfiable and reject otherwise.
\end{enumerate}
\end{definition}

A proof that $\succinct$ is $\NEXP$ complete can be found in~\cite[Chapter 20]{Pap94},
albeit with a different encoding.
Below, we show this implies $\NEXP$-completeness for our encoding as well.

\begin{proposition}
$\succinct$ is $\NEXP$-complete.
\end{proposition}
\begin{proof}
Papadimitriou~\cite{Pap94}
considers circuits~$\calC_{\mathrm{Pap}}$ which encode $\mathsf{3Sat}$ formulas~$\phi$ with $n$ variables and~$m$ clauses
as follows:
$\calC_{\mathrm{Pap}}$ takes as input a string $(b, u, k)$, where $b, k \in \{0, 1\}^2$ are interpreted as integers in $\{0, 1, 2, 3\}$
and $u \in \{0, 1\}^{\log(m)}$ is interpreted either as a vertex $1 \leq u \leq n$ or a clause $1 \leq u \leq m$.
If $1 \leq u \leq n$ and $0 \leq k \leq 2$, 
then on input $(0, u, k)$, $\calC_{\mathrm{Pap}}$ outputs the index of the clause where $\overline{x_u}$ appears for the $k$-th time,
and on input $(1, u, k)$,  it outputs the index of the clause where $x_u$ appears for the $k$-th time.
(In addition, if $1 \leq u \leq m$ and $0 \leq k \leq 3$, then on input $(2, u, k)$,
$\calC_{\mathrm{Pap}}$ outputs the $k$-th literal of the $u$-th clause in $\phi$.
We state this for completeness, though we will not need it for the proof.)
Such a $\mathsf{3Sat}$ formula~$\psi$ has $2n$ literals, each occurring~$3$ times, and so $m = 2n$.
By~\cite[Chapter 20]{Pap94}, this succinct encoding of $\mathsf{3Sat}$ is $\neexp$-complete.
Using this, we can generate an instance of the $\succinct$ problem~$\calC$ such that $\phi_\calC = \phi$ as follows:
given input $(u_1, u_2, u_3, b_1, b_2, b_3)$, we simply evaluate $\calC_{\mathrm{Pap}}$ on $(b_i, u_i, k)$, for each $1 \leq i \leq 3$ and $0 \leq k \leq 2$, and output~$1$ if there is any clause containing all three literals.
\end{proof}

The complete problem for $\NEEXP$ is, appropriately enough, a succinct
version of $\succinct$. To define it precisely, it helps to fix a
notion of a Boolean circuit. Following Section~4.3 of~\cite{Pap94}, we consider Boolean circuits in which each gate can
be one of six types: $\mathsf{input}, \mathsf{true}$, $\mathsf{false}$, $\wedge$,
$\vee$, or $\lnot$. These gates have 0, 0, 0, 2, 2, or 1 inputs,
respectively. A succinct representation of a circuit $\calC_1$ is a
circuit $\calC_2$ that, given an index $i$, outputs the type of gate
$i$ as well as the indices $j_1, j_2$ of its inputs (one or both of
these indices may be the null index $\varnothing$ depending on the
type of the gate $i$).
\begin{definition}
  $\succinctsquared$ is the following problem.
  \begin{itemize}
    \item[$\circ$] {\bf Input:} a circuit $\calC$ with size
      $\poly(n)$, which is a succinct representation of
      a circuit $\calC'$, which is itself an instance of $\succinct$
      with instance size $N = 2^{\poly(n)}$.
      \item[$\circ$] {\bf Output:} accept if $\psi_{\calC'}$ (the $\sat$
        formula on $2^N = 2^{2^{\poly(n)}}$ variables generated by the
        circuit $\calC'$) is
        satisfiable and reject otherwise.
      \end{itemize}
\end{definition}
\begin{fact}\label{fact:p-to-ppoly}
  Let $M$ be a deterministic Turing machine which takes two inputs $x_1, x_2$. Then
  for any input $x_1$ of size $n_1$ and for any size parameter $n_2$
  and time $T > n_1 + n_2$, there exists a
  circuit $C_{M,T,x_1}$ of size $N = O(T^2)$ which, on an input $x_2$ of
  size $n_2$, computes $M$ run for $T$
  steps on the input pair $x_1, x_2$. Moreover, there exists a Turing
  machine $M'$ that given $x_1$, $n_2$, and an
  index $i \in  \{1, \dots, N\}$ in binary, outputs in polynomial time
  the type of the $i$th gate of
  $C_{M,T, x_1}$ and the indices $j_1, j_2'$ of the inputs to this gate.
\end{fact}
\begin{proof}
  The construction in the proof of Theorem~8.1 of~\cite{Pap94} yields
  a circuit of the desired size. This circuit consists of $O(T^2)$
  copies of a constant-sized circuit $C_{M}$ that depends only on
  $M$. 
\end{proof}
\begin{theorem}[Cook-Levin]\label{thm:cook-levin}
  Let $L$ be a language in $\NTIME(T(n))$. Then the following
  properties hold:
  \begin{enumerate}
  \item For every string $x$ of length $n$, there exists a $\sat$
    formula $\Phi_x$ on $n' = \poly(T(n))$ variables $z_1, \dots, z_{n'}$, such that $x \in L$ iff
    $\Phi_x$ is satisfiable.
    \item There exists a Turing
      machine $R$ that given an input $x$ of
      length $n$, three indices $u_1, u_2, u_3 \in \{1, \dots, n'\}$ in
      binary, and three bits $b_1, b_2, b_3 \in \{0,1\}$, runs in
      $\poly\log(n') = \poly\log(T(n))$ time and outputs $1$ iff the clause
      $(z_{u_1}^{b_1}, z_{u_2}^{b_2}, z_{u_3}^{b_3})$ is included in $\Phi_x$.
    \end{enumerate}
  \end{theorem}
  \begin{proof}
    We follow the proof of Theorem~8.2 of~\cite{Pap94} to obtain the
    $\sat$ instance $\Phi_x$. 
  \end{proof}

  
\begin{theorem}
  $\succinctsquared$ is complete for $\NEEXP$ under polynomial time
  mapping reductions. That is, for any language $L$ in $\NEEXP$,
  there exists a Turing machine $R$ which takes as input a string $x
  \in \{0,1\}^n$, and in time $\poly(n)$ outputs an instance $\calC_{x}$ of
  $\succinctsquared$, such that $\calC_{x}$ is satisfiable iff $x \in L$.
\end{theorem}
\begin{proof}
  Suppose we start with a language $L \in \NEEXP$. This means there is
  a nondeterministic Turing machine $M$ which decides $L$ in time
  $T_0 = 2^{2^{n^c}}$.
  By \Cref{thm:cook-levin}, there exists a Turing machine $R_1$
  which runs in time $T_1 = \poly\log(T_0) = 2^{O(n^c)}$ such that
  given input $x$ and clause indices $i = (u_1, u_2, u_3, b_1, b_2, b_3)$,
  represented as a binary string of length $|i| = \log(\poly(T_0)) =
  \log(2^{O(2^{n^c})}) = O(2^{n^c})$, runs in time polynomial in
  $\ell$ and outputs $1$ iff the corresponding clause exists in a
  $\sat$ formula  $\Phi_x$ such that $x \in L$ iff $\Phi_x$ is satisfiable. 

  Now, if we apply~\Cref{fact:p-to-ppoly} to $R_1$, with $x$ playing
  the role of the first input $x_1$ and $i$ the role of the second
  input, we obtain that for every $x$ there exists a circuit $C_{R_1, T_1,
    x}$ of size $O(T_1^2) = 2^{O(n^c)}$ which takes as input a tuple
  of indices $i$, and runs $R_1$ for time $T_1$ on this time to output
  whether clause $i$ is present in the formula $\Phi_x$. Moreover, there exists a
  Turing machine $R_2$ that, given $x$, the size parameter $|i| =
  O(2^{n^c})$, represented in binary as a string of $O(n^c)$ bits, and an index $j$, represented
  as a string of $O(n^c)$ bits, outputs
  the $j$th gate of $C_{R_1, T_1, x}$ in time $T_2 = \poly(n)$. Note that
  $R_2$ is a Turing machine which takes in input of size $\poly(n)$
  and runs in time $\poly(n)$.

  We are now almost where we need to be. In the final step, we once
  again apply \Cref{fact:p-to-ppoly} to $R_2$, obtaining a third Turing machine $R_3$
  that takes as input $x$ and the size parameters, and an index $k$,
  and generates the $k$th gate of the circuit $C_{R_2, T_2, x}$
  corresponding to running $R_2$ for $T_2$ steps. Finally, by
  fixing the dependence of the size parameters on the size of $x$, and
  iterating through all possible values of the index parameter, we
  obtain a Turing machine $R_3'$ that takes as input $x$ and runs in
  time $\poly(n)$, and outputs the complete description of a $\succinctsquared$ instance $\calC_x$
  with the desired properties.
\end{proof}

\subsection{The Tseitin transformation}

In this section, we introduce the Tseitin transformation, which is a simple method of converting a Boolean circuit into a Boolean formula.

\begin{definition}[Tseitin transformation]
Let~$\calC$ be a Boolean circuit with~$n$ input variables $x_1, \ldots, x_n$ and~$s$ gates.
Then the \emph{Tseitin transformation of $\calC$}, denoted $\calF := \mathrm{Tseitin}(\calC)$, is the Boolean formula defined as follows.
\begin{itemize}
\item[(i)] Introduce new variables $w_1, \ldots, w_s$ corresponding to the output wires of the gates in~$\calC$.
		Then the input variables to $\calF$ consist of  $x_1, \ldots, x_n$ along with $w_1, \ldots, w_s$.
\item[(ii)] Each gate in~$\calC$ operates on one or two variables in $\{x_1, \ldots, x_n, w_1, \ldots, w_s\}$.
		Write $g_i(x, w)$ for the function computed by the $i$-th gate.
		Then $\calF$ computes the intermediate expression
		\begin{equation*}
			z_i := (g_i(x, w) \land w_i) \lor (\overline{g_i(x, w)} \land \overline{w_i}).
		\end{equation*}
		The final output of $\calF$ is $z_1 \land (z_2 \land ( \cdots \land z_s))$.
\end{itemize}
By construction, $\calC(x) = 1$ if and only if there exists a~$w$ such that $\calF(x, w) = 1$
(in particular, $w$ is taken to be the wire values of~$\calC$ on input~$x$).
In addition, $\calF$ contains exactly $7s + (s-1)$ gates, meaning that it has size~$O(s)$.
\end{definition}

Next, we show how to convert Boolean formulas into functions over $\F_q$.

\begin{definition}[Arithmetization]\label{def:arithmetization}
Let $\calF$ be a Boolean formula of~$n$ variables and size~$s$.
The \emph{arithmetization of $\calF$ over $\F_q$}, denoted $\arith{q}{\calF}$, is the formula produced by the following two-step process.
\begin{itemize}
\item[(i)] Transform $\calF$ by replacing all~$\lor$ gates with appropriate~$\land$ and~$\neg$ gates.
\item[(ii)] Transform each Boolean gate into an $\F_q$ gate as follows:
		Replace each~$\land$ gate in $\calF$ with a $\times$ gate.
		Replace each~$\neg$ gate with a $\times -1$ gate followed by a $+1$ gate (enacting the transformation $b \in \F_q \mapsto 1-b$).
		Call the resulting formula $\arith{q}{\calF}$.
\end{itemize}
Set $\calF_{\mathrm{arith}} := \arith{q}{\calF}$.
On inputs $x \in \{0, 1\}^n$, $\calF_{\mathrm{arith}}(x) = \calF(x)$.
On general inputs $x \in \F_q^n$, $\calF_{\mathrm{arith}}(x)$ is computable in time $\poly(s, q)$.
\end{definition}

The following proposition shows that small Boolean formulas have low-degree arithmetizations.

\begin{proposition}[Low-degree arithmetization]\label{prop:low-degree-arithmetization}
Let $\calF$ be a Boolean formula of~$n$ variables, size~$s$, and~$m$ gates.
Then $\arith{q}{\calF}$ is a degree-$s$ polynomial over~$\F_q$.
\end{proposition}
\begin{proof}
By induction on the number of gates, the base case $(m = 0)$ being trivial.
For the induction hypothesis, assume the proposition holds for Boolean formulas which have fewer than~$m$ gates.
Either the gate at the root of $\calF$ is a $\neg$ gate or an $\{\lor, \land\}$-gate.
In the former case, $\calF = \neg \calF'$ for some Boolean formula with $m-1$ gates,
and so $\arith{q}{\calF} = 1 - \arith{q}{\calF'}$ by construction.
But these have the same degree, and so $\arith{q}{\calF}$ is degree~$s$ by the induction hypothesis.
In the latter case, assume without loss of generality that it is an $\land$-gate. 
Then $\calF= \calF_{\mathrm{left}} \land \calF_{\mathrm{right}}$ for two formulas
of size $s_{\mathrm{left}} + s_{\mathrm{right}} = s$ and fewer than~$m$ gates.
By the induction hypothesis, $\arith{q}{\calF_{\mathrm{left}}}$ has degree-$s_{\mathrm{left}}$ and
$\arith{q}{\calF_{\mathrm{right}}}$ has degree-$s_{\mathrm{right}}$,
and so $\arith{q}{\calF} = \arith{q}{\calF_{\mathrm{left}}} \times \arith{q}{\calF_{\mathrm{right}}}$ has degree~$s$.
\end{proof}

The arithmetization procedure describe in \Cref{def:arithmetization} can also be applied to general Boolean circuits~$\calC$, not just Boolean formulas.
But \Cref{prop:low-degree-arithmetization} does not apply to general circuits;
in fact, the arithmetization of a Boolean circuit can have very high degree, even if that circuit is small.
This motivates using the Tseitin transformation: it allows us to convert a small circuit into a small formula, which has a low-degree arithmetization.


%%% Local Variables:
%%% mode: latex
%%% TeX-master: "main"
%%% End:
%%%---------- close: prelims-classical

%%%%%%%%%%% jump to prelims-quantum
%%%---------- open: prelims-quantum
%!TEX root = main.tex

